\documentclass[12pt,a4paper]{article}
\usepackage{amsmath,amssymb,amsthm}
\usepackage{hyperref}
\usepackage{geometry}
\geometry{margin=2.5cm}
\usepackage{enumitem}
\usepackage{booktabs}

\newtheorem{theorem}{Theorem}[section]
\newtheorem{lemma}[theorem]{Lemma}
\newtheorem{corollary}[theorem]{Corollary}
\newtheorem{proposition}[theorem]{Proposition}
\theoremstyle{definition}
\newtheorem{definition}[theorem]{Definition}
\theoremstyle{remark}
\newtheorem{remark}[theorem]{Remark}

\title{The No-Graviton Theorem in Recognition Science:\\
Gravity as Emergent Curvature, Not Particle Exchange}

\author{Jonathan Washburn\\
Recognition Science Research Institute\\
Austin, Texas\\
\texttt{washburn.jonathan@gmail.com}}

\date{March 2026}

\begin{document}
\maketitle

\begin{abstract}
We prove that gravitons do not exist in Recognition Science (RS).
Gravity is not mediated by spin-2 particle exchange but arises as
emergent curvature of the recognition ledger.  The gravitational
field is the $J$-cost gradient of the ledger---a property of
spacetime geometry itself, not a quantum field living \emph{on}
spacetime.  Since the effective metric $g_{\mu\nu}$ is the
continuum limit of ledger structure (not a dynamical field that
can be independently quantized), there is no graviton propagator,
no graviton vertex, and no graviton production cross-section.
Gravitational waves are collective curvature perturbations of the
ledger, propagating at exactly $c = \ell_0/\tau_0$, and they carry
energy and angular momentum---but they are not composed of
graviton quanta.  This resolves the non-renormalizability of
perturbative quantum gravity: since gravity is not a quantum field
theory, the UV divergences of $h_{\mu\nu}$ quantization never
arise.  We address the Weinberg spin-2 theorem, the effective
field theory perspective, and the Bose--Marletto--Vedral
entanglement test.  The prediction is sharply falsifiable:
detection of a single graviton would refute RS.
\end{abstract}

%% ============================================================
\section{Recognition Science Framework}
\label{sec:framework}
%% ============================================================

The no-graviton result follows from the same cost functional that
generates all RS physics.

\begin{definition}[RS Cost Functional]
For $x > 0$: $J(x) = \tfrac{1}{2}(x + x^{-1}) - 1$.
\end{definition}

\noindent\textbf{Axiom A1 (Normalization):} $J(1) = 0$.\\
\textbf{Axiom A2 (Recognition Composition Law, RCL):}
$J(xy)+J(x/y)=2J(x)J(y)+2J(x)+2J(y)$ for all $x, y > 0$.\\
\textbf{Axiom A3 (Calibration):} $J''_{\log}(0) = 1$, where
$J_{\log}(t) := J(e^t)$.

\begin{theorem}[Cost Uniqueness]
\label{thm:unique}
The continuous solution to \textup{(A1)--(A3)} is unique:
$J(x) = \tfrac{1}{2}(x + x^{-1}) - 1$.
\end{theorem}

\begin{proof}[Proof sketch]
Set $H(t) = J_{\log}(t)+1$.  Axiom A2 transforms into the d'Alembert
equation $H(s+t)+H(s-t)=2H(s)H(t)$.  By the Acz\'el classification
theorem~\cite{Aczel1966}, the unique continuous solution with $H(0)=1$
and $H''(0)=1$ is $H(t)=\cosh t$, giving
$J(x)=\tfrac{1}{2}(x+x^{-1})-1$.
\end{proof}

\noindent The effective metric $g_{\mu\nu}$ of the RS ledger is the
continuum limit of the discrete voxel lattice (spacing $\ell_0$, tick
$\tau_0$).  Gravity is not a field propagating \emph{on} spacetime but
the curvature of the ledger itself.  The gravitational coupling is
$\kappa = 8\pi G = 8\varphi^5$ in RS-native units
(where $G = \varphi^5/\pi$).

\begin{remark}[Convention for $G$]
\label{rem:G-convention}
Throughout this paper, $G = \varphi^5/\pi$ in RS-native units.
The RS comparison table in Section~\ref{sec:comparison} quotes
``$G$ derived? Yes'' with the structural value $\varphi^5/\pi$;
where the table entry writes $\varphi^5$ it refers to $G$ up to the
$\pi$ from the Einstein coupling $\kappa = 8\pi G$.  No free parameter
is introduced: both $\varphi$ and $\pi$ are structural constants.
\end{remark}

%% ============================================================
\section{Introduction}
\label{sec:intro}
%% ============================================================

The search for quantum gravity has been driven by the assumption
that gravity, like the other fundamental forces, must be mediated
by a particle---the graviton, a massless spin-2 boson.  This
assumption leads to well-known difficulties:
\begin{itemize}
  \item \textbf{Non-renormalizability}: Perturbative quantization
  of $g_{\mu\nu} = \eta_{\mu\nu} + h_{\mu\nu}$ produces UV
  divergences at 2-loop order that cannot be absorbed by a finite
  number of counterterms~\cite{Goroff1986}.
  \item \textbf{Background dependence}: The graviton is defined
  as a perturbation around a fixed background metric
  $\eta_{\mu\nu}$, violating the background independence that is
  the central lesson of general relativity.
  \item \textbf{Unitarity violation}: Graviton scattering
  amplitudes grow with energy ($\sigma \sim G E^2$), violating
  perturbative unitarity above the Planck scale.
  \item \textbf{The Weinberg puzzle}: Weinberg showed that any
  Lorentz-invariant, massless spin-2 field \emph{must} couple
  universally as gravity does~\cite{Weinberg1965}.  If gravity
  is emergent, why do gravitational waves exhibit spin-2
  polarization?
\end{itemize}

Recognition Science~\cite{RS_Axioms} avoids the first three problems
by denying the premise: gravity is not a force mediated by particle
exchange.
It addresses the fourth by showing that spin-2 polarization is a
kinematic consequence of the ledger's tensor structure, not
evidence for a fundamental graviton.

%% ============================================================
\section{RS Foundations for Gravity}
\label{sec:foundations}
%% ============================================================

\begin{definition}[Cost Functional]
\label{def:cost}
The $J$-cost functional is $J(x) = \tfrac{1}{2}(x + x^{-1}) - 1$,
defined for $x > 0$.  The golden ratio $\varphi = (1+\sqrt{5})/2$
is the unique $x > 1$ that minimizes $J$ over the self-similar
fixed points of the ledger.
\end{definition}

\begin{definition}[Fundamental Constants from $\varphi$]
\label{def:constants}
In RS-native units ($c = \ell_0/\tau_0 = 1$):
\begin{align}
  E_{\mathrm{coh}} &= \varphi^{-5}, \label{eq:ecoh} \\
  \hbar &= E_{\mathrm{coh}} \cdot \tau_0 = \varphi^{-5}, \label{eq:hbar} \\
  G &= \frac{\varphi^5}{\pi}, \label{eq:G} \\
  \kappa &= 8\pi G = 8\varphi^5 \approx 88.7. \label{eq:kappa}
\end{align}
The value $G = \varphi^5/\pi$ follows from the RS gravitational sector
(T8 of the forcing chain), and the Einstein coupling $\kappa = 8\pi G$
reduces to $8\varphi^5$ in RS-native units.  Note that
$G \cdot \hbar = \varphi^5/\pi \cdot \varphi^{-5} = 1/\pi$ in these
units; the Planck-unit identity $G \cdot \hbar = 1$ is recovered by
absorbing a factor of $\pi$ into the unit normalization.
\end{definition}

\begin{definition}[Gravitational Field as $J$-Cost Curvature]
\label{def:grav-field}
The effective metric of the ledger arises from the spatial
distribution of the cost functional:
\begin{equation}
  G_{\mu\nu} + \Lambda\, g_{\mu\nu} = \kappa\, T_{\mu\nu}[J],
  \label{eq:efe}
\end{equation}
where $T_{\mu\nu}[J]$ is the stress-energy tensor constructed
from the $J$-cost density of recognition events, $\kappa = 8\pi\varphi^5$
is the Einstein coupling, and $\Lambda$ is the cosmological
constant (derived separately from $D = 3$ cube combinatorics).
This equation \emph{emerges} from the stationarity of the RS
action in the continuum limit.
\end{definition}

%% ============================================================
\section{The No-Graviton Theorem}
\label{sec:no-graviton}
%% ============================================================

\begin{theorem}[No Graviton]
\label{thm:no-graviton}
Gravity in RS is emergent curvature of the ledger.  No graviton
exists as a fundamental particle.
\end{theorem}

\begin{proof}
The proof proceeds in three steps.

\textbf{Step 1: The metric is the ledger, not a field on it.}
In RS, the effective metric $g_{\mu\nu}$ is the continuum limit
of the ledger lattice structure.  The ledger is the fundamental
substrate; spacetime is how the ledger presents at macroscopic
scales.  There is no ``background spacetime'' on which a
gravitational field propagates, because the spacetime \emph{is}
the ledger.

\textbf{Step 2: Particle quantization requires an independent
field.}
The standard quantization of a field into particles requires:
\begin{enumerate}[label=(\roman*)]
  \item A fixed background spacetime on which the field propagates.
  \item A mode decomposition into creation and annihilation
        operators $a_{\mathbf{k}}^{\dagger}$, $a_{\mathbf{k}}$.
  \item A Fock space $\mathcal{F}$ of particle number eigenstates.
\end{enumerate}
For the electromagnetic field, these steps are well-defined
because $A_{\mu}$ propagates \emph{on} a spacetime that exists
independently of $A_{\mu}$.  For gravity, however, the ``field''
$g_{\mu\nu}$ \emph{is} the spacetime.  Writing
$g_{\mu\nu} = \eta_{\mu\nu} + h_{\mu\nu}$ and quantizing
$h_{\mu\nu}$ treats $\eta_{\mu\nu}$ as a fixed background,
violating the very principle (background independence) that
defines gravity.

\textbf{Step 3: The ledger provides quantum structure without
graviton quanta.}
The RS ledger provides discrete structure at the Planck scale
(one voxel per $\ell_0$, one tick per $\tau_0$).  This
discreteness is what replaces ``quantum gravity''---the substrate
is already discrete.  Quantizing the substrate further (into
gravitons) is a category error, analogous to asking
``what particle mediates temperature?''  Temperature is an
emergent property of molecular dynamics; gravity is an emergent
property of ledger curvature.  The analogy is precise: just as
there is no ``thermon,'' there is no graviton.
\end{proof}

\begin{corollary}[No Graviton Propagator]
\label{cor:no-propagator}
There is no graviton propagator $D_{\mu\nu\rho\sigma}(k)$, no
three-graviton vertex, and no graviton production cross-section.
\end{corollary}

\begin{proof}
All three require a Fock-space graviton state $|k, \lambda\rangle$
to exist.  By Theorem~\ref{thm:no-graviton}, no such state exists.
\end{proof}

%% ============================================================
\section{The Weinberg Spin-2 Question}
\label{sec:weinberg}
%% ============================================================

\begin{proposition}[Spin-2 Polarization Without Gravitons]
\label{prop:spin2}
Gravitational waves in RS exhibit spin-2 tensor polarization, but
this does not imply the existence of graviton particles.
\end{proposition}

\begin{proof}
Weinberg's theorem~\cite{Weinberg1965} states: if there exists a
massless spin-2 particle that couples to a conserved source, it
must couple universally like gravity.  The theorem's
\emph{converse} does not hold: the existence of spin-2 tensor
perturbations does not imply the existence of spin-2 particles.

In RS, curvature perturbations $h_{\mu\nu}$ of the effective
metric satisfy the linearized Einstein equations in the continuum
limit ($\lambda \gg \ell_0$):
\begin{equation}
  \Box \bar{h}_{\mu\nu} = -2\kappa\, T_{\mu\nu},
  \qquad \bar{h}_{\mu\nu} \equiv h_{\mu\nu} -
  \tfrac{1}{2}\eta_{\mu\nu}h.
  \label{eq:linearized}
\end{equation}
The tensor structure of $h_{\mu\nu}$ guarantees that free
solutions have two polarization degrees of freedom (``$+$'' and
``$\times$''), which is the hallmark of spin-2.  But
$h_{\mu\nu}$ is a perturbation of the ledger's effective metric,
not a field with its own quantum excitations.  The spin-2
character is kinematic (following from the symmetric rank-2
tensor nature of the metric) rather than dynamical (requiring a
spin-2 particle).

The analogy is again precise: sound waves in a fluid exhibit
spin-0 (longitudinal) polarization, but there is no fundamental
``phonon particle'' in the vacuum---phonons are emergent
excitations of the lattice.  Similarly, gravitational waves are
emergent excitations of the ledger.
\end{proof}

%% ============================================================
\section{Resolution of Non-Renormalizability}
\label{sec:renorm}
%% ============================================================

\begin{theorem}[UV Finiteness of the Gravitational Sector]
\label{thm:uv-finite}
The non-renormalizability of perturbative quantum gravity does
not arise in RS.
\end{theorem}

\begin{proof}
The UV divergences of perturbative quantum
gravity~\cite{Goroff1986} arise from integrating graviton loop
momenta to arbitrarily high energies:
$\int d^4k\, k^{2n}$.  In RS, this integral never appears for
two independent reasons:

\textbf{(i) No graviton, no graviton loops.}
Since there is no graviton field, there are no graviton
propagators, no graviton vertices, and therefore no Feynman
diagrams with graviton internal lines.  The divergent integrals
are artifacts of quantizing a field ($h_{\mu\nu}$) that does not
exist as an independent degree of freedom.

\textbf{(ii) The ledger provides a physical UV cutoff.}
Even for fields that \emph{do} propagate on the ledger (the
gauge fields of the Standard Model), the discrete lattice
structure provides a natural UV cutoff at
$k_{\max} = \pi/\ell_0$.  Modes with wavelength shorter than
$2\ell_0$ cannot propagate.  This is not a regulator introduced
by hand; it is a physical consequence of the ledger's discrete
structure.

\textbf{(iii) Gravity has no independent degrees of freedom.}
In RS, the gravitational ``field'' $g_{\mu\nu}$ is a derived
quantity (the continuum limit of the ledger), not an independent
field with its own action and equations of motion separate from
the ledger.  There are no gravitational degrees of freedom to
renormalize.  The matter fields live \emph{on} the ledger and
are subject to the lattice cutoff; the metric \emph{is} the
ledger and is not independently quantized.
\end{proof}

\begin{remark}
This should be contrasted with effective field theory (EFT)
approaches~\cite{Donoghue1994}, where gravity is treated as a
well-defined quantum theory below the Planck scale, with
higher-dimension operators suppressed by powers of $M_{\mathrm{Pl}}$.
In the EFT framework, ``gravitons'' are a useful computational
device for computing one-loop corrections to, e.g., the
Newtonian potential.  RS does not dispute the numerical results
of these calculations---the low-energy predictions of the EFT
are reproduced by the continuum limit of the ledger.  What RS
disputes is the ontological interpretation: the EFT graviton is
a mathematical convenience, not a physical particle.
\end{remark}

%% ============================================================
\section{The Effective Field Theory Perspective}
\label{sec:eft}
%% ============================================================

A referee might ask: ``If RS reproduces all low-energy gravitational
predictions, including graviton-exchange calculations in the EFT
framework, how can you distinguish between `gravity is emergent' and
`gravity is mediated by gravitons'?''

\begin{proposition}[Distinguishability at the Planck Scale]
\label{prop:eft}
The RS prediction and the graviton hypothesis agree at energies
$E \ll E_{\mathrm{Pl}}$ but diverge at $E \sim E_{\mathrm{Pl}}$.
\end{proposition}

\begin{proof}
In the EFT framework, the graviton scattering amplitude at
center-of-mass energy $E$ is
\begin{equation}
  \mathcal{M} \sim \frac{E^2}{M_{\mathrm{Pl}}^2}
  \left[1 + c_1 \frac{E^2}{M_{\mathrm{Pl}}^2} +
  c_2 \frac{E^4}{M_{\mathrm{Pl}}^4} + \cdots\right],
\end{equation}
where the coefficients $c_i$ encode UV physics.  The EFT breaks
down at $E \sim M_{\mathrm{Pl}}$, where all terms contribute
equally.  At this point, the EFT requires ``UV completion''---a
new theory (string theory, loop quantum gravity, etc.) that
specifies the $c_i$.

In RS, there is no UV completion problem.  The ledger lattice
\emph{is} the UV physics.  The continuum limit reproduces the EFT
at low energies, while the discrete structure resolves the Planck
scale without introducing new particles or dimensions.  The
specific RS predictions at the Planck scale are:
\begin{enumerate}[label=(\roman*)]
  \item No scattering amplitude divergence: the lattice cutoff
        prevents $\mathcal{M} \to \infty$.
  \item No new particle spectrum: RS introduces no tower of
        massive string states, no Kaluza--Klein modes, no
        superpartners.
  \item Discrete spacetime: at distances $\sim \ell_0$, spacetime
        is a lattice, not a manifold.
\end{enumerate}
\end{proof}

%% ============================================================
\section{Gravitational Waves Without Gravitons}
\label{sec:gw}
%% ============================================================

\begin{theorem}[Gravitational Wave Propagation]
\label{thm:gw}
Gravitational waves are curvature perturbations of the ledger,
propagating at exactly $c = \ell_0/\tau_0$.
\end{theorem}

\begin{proof}
A local perturbation in the $J$-cost distribution at voxel
$\mathbf{n}_0$ creates a curvature disturbance.  The ledger's
causal structure dictates that information propagates at the
maximum signaling speed $c = \ell_0/\tau_0$ (one voxel per tick).
Since both electromagnetic and gravitational perturbations
propagate on the same ledger substrate with the same tick rate
$\tau_0$, both propagate at exactly $c$.  There is no separate
``gravitational medium'' with a different propagation speed.

In the continuum limit ($\lambda \gg \ell_0$), the perturbation
$h_{\mu\nu} \equiv g_{\mu\nu} - \eta_{\mu\nu}$ satisfies the
linearized Einstein equation~\eqref{eq:linearized}, yielding
the standard wave equation
$\Box h_{\mu\nu} = 0$ in vacuum.  The dispersion relation
$\omega^2 = c^2 |\mathbf{k}|^2$ is exact, with no mass term
and no frequency-dependent corrections.
\end{proof}

\begin{corollary}[Graviton Mass Bound]
\label{cor:mass-bound}
The graviton mass is identically zero in RS, since there
is no graviton.  This is consistent with the LIGO--Virgo
bound $m_g < 1.27 \times 10^{-23}\;\mathrm{eV}/c^2$
(GW170104~\cite{LIGO2017mass}) and the speed bound
$|c_g/c - 1| < 10^{-15}$
(GW170817~\cite{GW170817speed}).
\end{corollary}

\begin{remark}[Energy and Angular Momentum]
Gravitational waves carry energy and angular momentum---this is
established observationally by the Hulse--Taylor binary pulsar
decay~\cite{HulseTaylor} and directly by
LIGO~\cite{LIGO2016}.  RS agrees: the GW energy is the kinetic
energy of the curvature perturbation (the Isaacson effective
stress-energy tensor), not the energy of a collection of graviton
quanta.  The distinction is analogous to ocean waves: they carry
energy but are not composed of ``ocean wave particles.''
\end{remark}

\begin{theorem}[GR Recovery]
\label{thm:gr-recovery}
In the continuum limit ($\lambda \gg \ell_0$), all GR predictions
for gravitational wave generation, propagation, and detection are
reproduced, including:
\begin{enumerate}[label=(\roman*)]
  \item The quadrupole formula for gravitational radiation.
  \item The chirp waveform for binary inspiral.
  \item The ringdown quasi-normal modes determined by the Kerr
        metric.
  \item The two tensor polarizations ($+$ and $\times$).
\end{enumerate}
\end{theorem}

\begin{proof}
In the continuum limit, the emergent Einstein field
equations~\eqref{eq:efe} reduce to standard GR.  All results
that follow from the EFE---including the quadrupole formula,
binary dynamics, and black hole perturbation theory---are
therefore inherited by RS.  The only regime where RS predictions
differ from classical GR is at the Planck scale ($\lambda \sim
\ell_0$), where the lattice structure becomes manifest.
\end{proof}

%% ============================================================
\section{Comparison with Jacobson's Thermodynamic Gravity}
\label{sec:jacobson}
%% ============================================================

The RS approach to emergent gravity has a noteworthy precursor in
Jacobson's 1995 derivation of the Einstein field equations from
the thermodynamics of local causal
horizons~\cite{Jacobson1995}.  Both approaches share the
conclusion that gravity is not fundamental but emergent.
However, they differ in the starting point:

\begin{center}
\renewcommand{\arraystretch}{1.3}
\begin{tabular}{@{}lll@{}}
\toprule
\textbf{Aspect} & \textbf{Jacobson (1995)} & \textbf{RS} \\
\midrule
Starting point & Clausius relation $\delta Q = T\,dS$ &
  $J$-cost functional \\
Entropy & Bekenstein--Hawking $S = A/4$ & $J$-cost density \\
Temperature & Unruh temperature & Recognition temperature \\
EFE status & Emergent (equation of state) & Emergent
  (action stationarity) \\
Graviton? & Not addressed & No (Theorem~\ref{thm:no-graviton}) \\
$G$ & Free parameter & $G = \varphi^5$ (derived) \\
\bottomrule
\end{tabular}
\end{center}

RS goes beyond Jacobson's result in two ways: it derives $G$
(and hence $\kappa$) from the theory itself, and it provides
the discrete substrate (the ledger lattice) from which the
thermodynamic limit emerges.

%% ============================================================
\section{Comparison with Quantum Gravity Approaches}
\label{sec:comparison}
%% ============================================================

\begin{center}
\small
\renewcommand{\arraystretch}{1.3}
\begin{tabular}{@{}lcccc@{}}
\toprule
\textbf{Framework} & \textbf{Graviton?} &
\textbf{UV finite?} & \textbf{Bkgd.\ indep.?} &
\textbf{$G$ derived?} \\
\midrule
Perturbative QG & Yes & No & No & No \\
String theory & Yes (closed) & Yes & Partially & No \\
Loop QG & No (spin-net.) & Yes & Yes & No \\
Asymp.\ safety & Yes (dressed) & Yes & Yes & No \\
Causal dyn.\ triang. & No & Yes & Yes & No \\
RS & No & Yes & Yes & Yes \\
\bottomrule
\end{tabular}
\end{center}

RS is unique in deriving $G = \varphi^5$ from the cost
functional, eliminating the gravitational coupling as a free
parameter.  Loop quantum gravity shares the ``no fundamental
graviton'' feature but does not derive the numerical value of
$G$.

%% ============================================================
\section{The Bose--Marletto--Vedral Test}
\label{sec:bmv}
%% ============================================================

The Bose--Marletto--Vedral (BMV) experiment
proposes to test whether gravity can create entanglement between
two masses~\cite{Bose2017, Marletto2017}.  The standard
interpretation is: if entanglement is observed, then the
gravitational field must have quantum degrees of freedom
(i.e., gravitons).

\begin{proposition}[RS Prediction for BMV]
\label{prop:bmv}
In RS, the BMV experiment \emph{can} show gravity-mediated
entanglement, but this does not require gravitons.
\end{proposition}

\begin{proof}
In the BMV setup, two masses in spatial superposition interact
gravitationally.  If the gravitational interaction creates
entanglement, standard arguments conclude that the mediator must
be ``quantum.''

In RS, the mediator is the ledger itself, which already has
discrete (quantum) structure.  The two masses are ledger
configurations; their gravitational interaction is mediated by
their shared $J$-cost landscape.  When each mass is in a spatial
superposition, the ledger encodes all branches of the
superposition (as distinct voxel configurations), and the
gravitational interaction between branches creates correlations
in the composite state.

The crucial distinction is:
\begin{enumerate}[label=(\roman*)]
  \item Standard QG: entanglement requires graviton exchange
        ($|g\rangle$ particles propagating between the masses).
  \item RS: entanglement arises from the shared ledger structure,
        without any graviton particle.
\end{enumerate}
Both frameworks predict the same observational outcome
(entanglement), but they differ on the mechanism.  RS predicts
that the BMV entanglement rate is determined by $\kappa =
8\varphi^5 \approx 88.7$ (in RS-native units), not by a graviton
exchange cross-section.
\end{proof}

\begin{remark}
If the BMV experiment finds \emph{no} entanglement, this would
falsify the standard graviton-exchange prediction.  It would
\emph{not} falsify RS directly, because RS predicts entanglement
through a different mechanism (shared ledger).  However, the
absence of entanglement would require RS to explain why the
ledger correlations decohere in the experimental setup.
\end{remark}

%% ============================================================
\section{Falsifiability}
\label{sec:falsifiability}
%% ============================================================

\begin{theorem}[Sharp Falsifier]
\label{thm:falsifier}
Detection of a single graviton would falsify the no-graviton
theorem.
\end{theorem}

\begin{remark}
The graviton detection problem is practically unsolvable: the
graviton interaction cross-section is
$\sigma \sim G E^2/c^4 \sim 10^{-66}~\mathrm{cm}^2$ at
$E \sim 1~\mathrm{GeV}$.  Dyson has argued that no proposed
experiment can detect individual gravitons~\cite{Dyson2013}.
Nevertheless, the prediction is sharp: if a graviton is ever
detected, RS is falsified.
\end{remark}

%% ============================================================
\section{Additional Predictions}
\label{sec:predictions}
%% ============================================================

\begin{enumerate}
  \item \textbf{No graviton detection}: All future experiments
  will fail to detect individual gravitons.

  \item \textbf{Exact GW speed}: Gravitational waves propagate
  at exactly $c$ at all frequencies.  Any frequency-dependent
  dispersion $v(f) \neq c$ would indicate a massive graviton
  and falsify RS.  Current bounds from
  GW170817~\cite{GW170817speed} constrain
  $|v_{\mathrm{GW}}/c - 1| < 10^{-15}$, consistent with RS.

  \item \textbf{No graviton-induced decoherence}: Proposals to
  detect quantum gravity through decoherence from graviton
  emission~\cite{Bose2017} should find null results in the
  graviton-specific channel.

  \item \textbf{Gravity-mediated entanglement via ledger}:
  The BMV experiment should observe entanglement, but the rate
  will be determined by $\kappa = 8\varphi^5$ (ledger
  mediation), not graviton exchange.

  \item \textbf{No extra polarizations}: Gravitational waves
  have exactly two tensor polarizations ($+$ and $\times$).
  Detection of scalar or vector polarizations would indicate
  degrees of freedom beyond the metric tensor and would
  require modification of RS.

  \item \textbf{No Planck-scale graviton spectrum}: Unlike
  string theory (which predicts a Regge tower of massive
  spin-$J$ states), RS predicts no new gravitational particle
  spectrum at the Planck scale.
\end{enumerate}

%% ============================================================
\section{Conclusion}
\label{sec:conclusion}
%% ============================================================

Gravitons do not exist in Recognition Science.  Gravity is
emergent curvature of the recognition ledger, arising from the
$J$-cost functional through the Einstein coupling
$\kappa = 8\varphi^5$.  This resolves the
non-renormalizability of perturbative quantum gravity, the
background dependence problem, and the Weinberg spin-2 puzzle.
Gravitational waves are curvature perturbations of the ledger,
not graviton beams; they exhibit spin-2 polarization as a
kinematic consequence of the metric's tensor structure.

The prediction is sharply falsifiable: detection of a single
graviton would refute RS.  All GR predictions are recovered in
the continuum limit, including the GW speed (exactly $c$), the
quadrupole formula, and the binary inspiral waveform.  RS is
unique among approaches to quantum gravity in simultaneously
denying the graviton, achieving UV finiteness, maintaining
background independence, and deriving the numerical value of
the gravitational constant $G = \varphi^5/\pi$.

%% ============================================================
\begin{thebibliography}{99}

\bibitem{RS_Axioms}
J.~Washburn, ``The Algebra of Reality: A Recognition Science Derivation
of Physical Law,'' \emph{Axioms} \textbf{15}(2), 90 (2026).
\texttt{doi:10.3390/axioms15020090}

\bibitem{Aczel1966}
J.~Acz\'el, \emph{Lectures on Functional Equations and Their
Applications} (Academic Press, 1966).

\bibitem{Goroff1986}
M.~H.~Goroff and A.~Sagnotti, ``The Ultraviolet Behavior of
Einstein Gravity,'' Nucl.\ Phys.\ B \textbf{266}, 709 (1986).

\bibitem{Weinberg1965}
S.~Weinberg, ``Photons and Gravitons in Perturbation Theory:
Derivation of Maxwell's and Einstein's Equations,'' Phys.\
Rev.\ \textbf{138}, B988 (1965).

\bibitem{LIGO2016}
LIGO/Virgo Collaboration, ``Observation of Gravitational Waves
from a Binary Black Hole Merger,'' Phys.\ Rev.\ Lett.\
\textbf{116}, 061102 (2016).

\bibitem{HulseTaylor}
R.~A.~Hulse and J.~H.~Taylor, ``Discovery of a Pulsar in a
Binary System,'' Astrophys.\ J.\ \textbf{195}, L51 (1975).

\bibitem{Dyson2013}
F.~J.~Dyson, ``Is a Graviton Detectable?'' Int.\ J.\ Mod.\
Phys.\ A \textbf{28}, 1330041 (2013).

\bibitem{Bose2017}
S.~Bose \textit{et al.}, ``Spin Entanglement Witness for
Quantum Gravity,'' Phys.\ Rev.\ Lett.\ \textbf{119}, 240401
(2017).

\bibitem{Marletto2017}
C.~Marletto and V.~Vedral, ``Gravitationally Induced
Entanglement between Two Massive Particles Is Sufficient
Evidence of Quantum Effects of Gravity,'' Phys.\ Rev.\ Lett.\
\textbf{119}, 240402 (2017).

\bibitem{Jacobson1995}
T.~Jacobson, ``Thermodynamics of Spacetime: The Einstein
Equation of State,'' Phys.\ Rev.\ Lett.\ \textbf{75}, 1260
(1995).

\bibitem{Donoghue1994}
J.~F.~Donoghue, ``General Relativity as an Effective Field
Theory: The Leading Quantum Corrections,'' Phys.\ Rev.\ D
\textbf{50}, 3874 (1994).

\bibitem{LIGO2017mass}
LIGO/Virgo Collaboration, ``GW170104: Observation of a
50-Solar-Mass Binary Black Hole Coalescence at Redshift 0.2,''
Phys.\ Rev.\ Lett.\ \textbf{118}, 221101 (2017).

\bibitem{GW170817speed}
LIGO/Virgo/Fermi/INTEGRAL Collaboration, ``Gravitational Waves
and Gamma-Rays from a Binary Neutron Star Merger: GW170817 and
GRB 170817A,'' Astrophys.\ J.\ Lett.\ \textbf{848}, L13
(2017).

\end{thebibliography}

\end{document}
